\section{Results I: Parameter Count}
\label{sec:res-params}

Table~\ref{tab:params} breaks the parameter count of each configuration into its quantum and classical share.

\begin{table}[h]
\centering
\caption{Parameter breakdown. Quantum share is in percent of the total.}
\label{tab:params}
\begin{tabular}{@{}lrrrrr@{}}
\toprule
Config & Quantum & Classical & Total & Quantum share (\%) & Reduction vs.\ classical (\%) \\
\midrule
classical & 0  & 4353 & 4353 & 0.00 & --- \\
narrow    & 16 & 2405 & 2421 & 0.66 & 44.4 \\
baseline  & 32 & 2537 & 2569 & 1.25 & 41.0 \\
wide      & 48 & 2669 & 2717 & 1.77 & 37.6 \\
deep      & 64 & 2537 & 2601 & 2.46 & 40.2 \\
\bottomrule
\end{tabular}
\end{table}

The quantum branch contributes only 16--64 parameters (0.7--2.5\% of the total). Nearly all of the reduction
relative to the classical baseline therefore comes from the smaller classical part of the hybrid network
(two $32\times32$ hidden layers instead of five), not from the circuit.
\todo{Own paragraph: what this implies for the claim that quantum circuits are parameter-efficient.
A fair control would be a purely classical network of the same size; say whether it is in scope.}
